%!TEX root = ../Proeycto.tex
\chapter{Marco Teórico}

En este capítulo se desarrollan los fundamentos conceptuales y tecnológicos que respaldan el diseño, arquitectura e implementación del Sistema de Expedientes de Trámites y Citas (SETC). Se abordan los principios de la arquitectura Modelo-Vista-Controlador (MVC), el ecosistema de desarrollo en Laravel 13 y PHP 8.4, los modelos de seguridad basados en roles (RBAC), la asistencia interactiva mediante Inteligencia Artificial, la gestión documental digital y los esquemas de validación mediante sellos digitales y códigos QR.

\section{Arquitectura de Software y Patrón Modelo-Vista-Controlador}

El desarrollo de aplicaciones web empresariales requiere una separación clara de responsabilidades para garantizar mantenibilidad, modularidad y escalabilidad. El patrón arquitectónico Modelo-Vista-Controlador (MVC) desacopla la lógica de negocio, la presentación visual y el control del flujo de datos:

\begin{itemize}
    \item \textbf{Modelo:} Encapsula el estado de la aplicación, las reglas de validación de dominio, la integridad de los datos y las relaciones cardinales entre entidades.
    \item \textbf{Vista:} Representa la interfaz gráfica de usuario encargada de renderizar la información de manera accesible y responsiva.
    \item \textbf{Controlador:} Intercepta las solicitudes HTTP provenientes de los clientes, invoca la lógica de negocio o servicios correspondientes y retorna la respuesta adecuada (HTML, JSON o flujos binarios de documentos).
\end{itemize}

\section{Ecosistema Laravel 13 y PHP 8.4}

El sistema SETC está desarrollado sobre el lenguaje de programación \textbf{PHP 8.4}, aprovechando características avanzadas de tipado estricto, promoción de propiedades en constructores, enumeraciones tipadas y optimizaciones en el compilador JIT (\textit{Just-In-Time}), lo que se traduce en una reducción significativa del consumo de memoria y un menor tiempo de procesamiento por solicitud.

Sobre esta base, se adoptó el framework \textbf{Laravel 13}, el cual aporta componentes de ingeniería robustos:

\begin{itemize}
    \item \textbf{Eloquent ORM:} Mapeo objeto-relacional que permite interactuar con la base de datos utilizando sintaxis orientada a objetos expresiva, asegurando protección nativa contra inyecciones SQL mediante consultas parametrizadas PDO.
    \item \textbf{Motor de Plantillas Blade y Alpine.js:} Blade proporciona una estructura jerárquica de vistas con componentes reutilizables y directivas de seguridad contra ataques XSS (\textit{Cross-Site Scripting}). Por su parte, \textbf{Alpine.js} brinda reactividad declarativa y liviana en el navegador del cliente para el manejo de estados dinámicos, validación de formularios en tiempo real y modales interactivos.
    \item \textbf{Tailwind CSS:} Framework CSS utilitario que permite construir interfaces responsivas, accesibles y consistentes bajo un diseño ergonómico de alta productividad.
    \item \textbf{Canalización de Middlewares:} Permite encadenar capas de inspección previas a la ejecución de los controladores para garantizar la autenticación, autorización y registro de bitácora.
\end{itemize}

\section{Seguridad de la Información y Control de Acceso Basado en Roles (RBAC)}

La protección de expedientes telemáticos gubernamentales exige políticas de control de acceso estrictas. El modelo de \textbf{Control de Acceso Basado en Roles} (\textit{Role-Based Access Control}, RBAC) asigna permisos a roles definidos en lugar de a usuarios individuales, simplificando la administración y previniendo brechas de privilegios.

En el sistema SETC, se modelaron cuatro roles institucionales:
\begin{enumerate}
    \item \textbf{Ciudadano:} Acceso restringido al catálogo público, inicio de solicitudes propias, consulta de su propia línea de tiempo, subsanación de recaudos observados y descarga de resoluciones y citas.
    \item \textbf{Revisor Técnico:} Operación de la Mesa de Dictamen para evaluar expedientes asignados, revisar recaudos digitales, emitir observaciones y formalizar dictámenes técnicos.
    \item \textbf{Supervisor / Dictaminador:} Supervisión global de expedientes y coordinación del flujo pericial en la consola de distribución de trabajo.
    \item \textbf{Administrador:} Gestión integral de catálogos de trámites, requisitos, usuarios, roles institucionales y auditoría del sistema.
\end{enumerate}

\section{Asistentes Basados en Inteligencia Artificial en Entornos Telemáticos}

La integración de asistentes conversacionales dotados de Inteligencia Artificial en plataformas de gobierno electrónico representa una evolución sustancial en la experiencia del usuario. Estos agentes permiten procesar las consultas en lenguaje natural de los ciudadanos, guiándolos paso a paso a través de los requisitos documentales exigibles para cada trámite específico. Al operar de forma interactiva y contextual, el asistente reduce el índice de errores en la radicación inicial de documentos, disminuye la necesidad de subsanaciones posteriores y alivia la saturación en los canales convencionales de soporte institucional.

\section{Gestión Documental Digital y Validación Criptográfica}

\subsection{Almacenamiento Estructurado y Formato PDF}
El estándar \texttt{PDF} (\textit{Portable Document Format}) garantiza que los recaudos adjuntados por la ciudadanía y las resoluciones emitidas por la autoridad mantengan su integridad visual, metadatos y fidelidad independientemente del dispositivo en que se visualicen. El sistema implementa almacenamiento estructurado con asignación de nombres basados en hashes criptográficos (\texttt{SHA-256}) y validación de cabeceras MIME reales en el servidor para impedir la carga de archivos no autorizados.

\subsection{Validación Telemática mediante Códigos QR}
Para permitir la verificación inmediata de la autenticidad de las resoluciones administrativas y pases digitales de cita, el sistema integra la generación dinámica de \textbf{Códigos QR} (\textit{Quick Response Code}). Estos códigos codifican un identificador unívoco y firma digital que enlaza al registro oficial centralizado, permitiendo a cualquier autoridad o verificador comprobar al instante la validez del documento mediante un dispositivo móvil.
